\documentclass[11pt]{article}
\usepackage[a4paper,margin=27mm]{geometry}
\usepackage{amsmath,amssymb,amsthm,xcolor,booktabs,array}
\newcommand{\PROVED}{\textcolor{green!40!black}{\textbf{PROVED}}}
\newcommand{\OPEN}{\textcolor{red!70!black}{\textbf{OPEN}}}
\newcommand{\AUDIT}{\textcolor{orange!70!black}{\textbf{AUDIT}}}
\newtheorem{theorem}{Theorem}
\newtheorem{corollary}{Corollary}
\title{The functional-pair metric problem for the odd Witt action (v107)}
\date{13 August 2026}

\begin{document}
\maketitle

\section{Allowing the most general positive coupling}

Let $\rho$ be a spectral character in the odd row-(c) object and pair it
with its Fourier--Poisson mate $1-\overline\rho$.  On the semisimple paired
block, the half-Tate centered dilation has matrix
\begin{equation}
 A_{n,\rho}=
 \begin{pmatrix}
 n^{1/2-\rho}&0\\
 0&n^{\overline\rho-1/2}
 \end{pmatrix}.                                                \tag{1}
\end{equation}
In particular $|\det A_{n,\rho}|=1$.  We do not assume that the two
spectral lines are orthogonal.  Let $P$ be an arbitrary positive definite
$2\times2$ Hermitian matrix and use
$\langle v,w\rangle_P=w^*Pv$.

\begin{theorem}[Determinant rigidity of a functional pair --- \PROVED]
For one integer $n\ge2$, the inequality
\begin{equation}
                       A_{n,\rho}^*PA_{n,\rho}\le P            \tag{2}
\end{equation}
holds for some positive definite $P$ if and only if
\begin{equation}
                         \Re\rho=\frac12.                       \tag{3}
\end{equation}
Whenever (2) holds, equality holds in (2): the paired action is $P$-unitary.
\end{theorem}

\begin{proof}
Put $T=P^{1/2}A_{n,\rho}P^{-1/2}$.  Inequality (2) is exactly
$T^*T\le I$.  Hence both singular values $s_1,s_2$ of $T$ are at most one.
Similarity and (1) give
\[
 s_1s_2=|\det T|=|\det A_{n,\rho}|=1.
\]
Therefore $s_1=s_2=1$, so $T^*T=I$ and equality holds in (2).  Every
eigenvalue of a unitary matrix has modulus one.  The eigenvalues in (1)
have moduli
\[
 n^{1/2-\Re\rho},\qquad n^{\Re\rho-1/2};
\]
both are one precisely when (3) holds.  Conversely, if (3) holds, (1) is
unitary for $P=I$.
\end{proof}

This theorem rules out a possible hidden closure: off-diagonal entries in
the Hilbert metric cannot absorb the expanding member of a functional pair.
The product of the two expansion factors is one, so a contraction of the
pair is automatically a unitary action.

\section{Jordan blocks and multiplicity}

Let a finite Loewy block at $\rho$ have the form
\begin{equation}
 A=n^{1/2-\rho}\exp(-N\log n),                                 \tag{4}
\end{equation}
where $N$ is nilpotent.  Pair it with its functional mate so that the
determinant of the total centered action again has modulus one.

\begin{corollary}[Contractivity kills the nilpotent part --- \PROVED]
If the paired finite Loewy block is contractive for a positive definite
metric, then $\Re\rho=1/2$ and $N=0$ on its semisimplified unitary part.
In particular a trace-preserving faithful Hilbert realization of the full
Loewy block would force both the critical line and semisimplicity.
\end{corollary}

\begin{proof}
The determinant argument in Theorem 1 makes the whole paired operator
unitary.  A unitary matrix is normal and hence diagonalizable.  Its
restriction to each generalized eigenspace therefore has no nontrivial
nilpotent Jordan part.  The eigenvalue-modulus argument gives
$\Re\rho=1/2$.
\end{proof}

\section{Exact consequence for the proposed completion}

The preceding result uses every positive metric on the paired block; it is
not tied to the diagonal Witt norm, to ordinary $L^2$, or to a choice of
basis.  Consequently a construction from rows (a)--(c) of a positive
trace-faithful completion with the Witt adjoint would itself prove
\begin{equation}
 \Re\rho=\frac12\quad\text{for every odd spectral character }\rho. \tag{5}
\end{equation}
There is no remaining finite-dimensional metric parameter to tune.  The
new source theorem must instead prove positivity of the canonical
Fourier--Poisson pairing on every functional pair.  In block form this is
the assertion that the canonical Hermitian Gram $P_\rho$ is positive;
Theorem 1 would then yield (5) without circularity because $P_\rho$ would
have been constructed and proved positive before spectral decomposition.

\begin{center}
\begin{tabular}{>{\raggedright\arraybackslash}p{.58\linewidth}
                >{\centering\arraybackslash}p{.15\linewidth}
                >{\raggedright\arraybackslash}p{.17\linewidth}}
\toprule
Statement & Status & Location \\
\midrule
Arbitrary positive metric on a functional pair & classified & Theorem 1 \\
Pair contractivity forces criticality and unitarity & \PROVED & (1)--(3) \\
Faithful Loewy completion forces semisimplicity & \PROVED & Corollary 2 \\
Source construction and positivity of the canonical pairing $P_\rho$ & \OPEN & not supplied by scalar contact \\
Condition $G$ and row (d) & \OPEN & follow from the preceding source theorem \\
\bottomrule
\end{tabular}
\end{center}

\end{document}
